\documentclass[runningheads]{llncs}

\usepackage{amsmath,amssymb,mathtools}
\usepackage{array}
\usepackage{booktabs}
\usepackage{enumitem}
\usepackage[hidelinks]{hyperref}

\newcommand{\OPT}{\mathrm{OPT}}
\newcommand{\PW}{\mathrm{PW}}
\newcommand{\FAIR}{\mathrm{FAIR}}
\newcommand{\PoF}{\mathrm{PoF}}
\newcommand{\E}{\mathbb{E}}
\newcommand{\R}{\mathbb{R}}
\newcommand{\argmax}{\operatorname*{arg\,max}}
\newcommand{\argmin}{\operatorname*{arg\,min}}

\begin{document}

\title{The Price of Fairness, Complexity, and Incentives in
No-Numeraire Combinatorial Trade-Intent Auctions}
\titlerunning{No-Numeraire Combinatorial Trade-Intent Auctions}

\author{Anonymous Author(s)}
\authorrunning{Anonymous}
\institute{Anonymous Institution}

\maketitle

\begin{abstract}
Blockchain trade-intent auctions, which currently intermediate billions
of dollars in trades per month, are combinatorial---a solver executing
several intents jointly can net opposing flows and save on fees---yet
they have no numeraire: there is no common asset in which to redistribute
the resulting surplus. The deployed fix builds an endogenous fairness
benchmark from the bids and executes a batched bid only if it delivers to
every covered trader at least that trader's standalone competitive
outcome. We give a worst-case theory of this design using a common scale
parameter: the complementarity factor $g$, together with liquidity
$\theta=1-\lambda$.
The model and fairness benchmark are due to Canidio and
Henneke~\cite{canidio2024}; our contribution is the worst-case
approximation, competitive, and incentive layer on top.

The price of fairness is exactly $\min(g,1/\theta)$, tight for all
parameters and independent of the number of orders and of prices, with a
dichotomy between $g$ (no numeraire) and $1$ (full numeraire); batch size
is not the right parameter---the right refinement is a per-leg coverage
floor $\alpha$, under which the exact law is
$\min(g,\alpha+(1-\alpha)/\theta)$. Within the order-price abstraction of
UDP, the rate-realizable floor-vector frontier for equilibrium support of
efficient allocations is exactly $g=1$, by native no-numeraire arguments;
at $g=1$ the price system is precisely the fairness benchmark. The price of
strategyproofness is $\Theta(1+\log g)$ for a single multi-minded solver
and for many single-minded solvers, via two posted-price mechanisms and a
convex profit-curve technique. The remaining multi-agent multi-minded
cell is a partial frontier: it admits $O(N\log g)$ and has a serial-demand
$\Omega(\min(N,g))$ barrier. Online, the price splits: against realized
lifetime floors even randomization pays exactly $g$, while with static
public floors the residual problem is one-way search over the spread,
deterministic $\Theta(\sqrt g)$ and randomized $\Theta(\log g)$.
Across these results, the same scale-free parameter $g$ recurs in both
the informational and the incentive price of the spread.
\keywords{Combinatorial auctions \and Price of fairness \and
Mechanism design \and Blockchain \and Competitive analysis.}
\end{abstract}

\section{Introduction}\label{sec:intro}

Blockchain ``trade-intent'' auctions let users state a desired exchange
of one asset for another and delegate execution to competing agents
called \emph{solvers}; the protocol collects intents off-chain and
auctions them in batches, in the spirit of frequent-batch market
design~\cite{budish2015}. These markets are economically significant---
recent figures put monthly volume on the order of billions of
dollars\footnote{Order-of-magnitude figure from the anonymized deployment
diagnostics in Appendix~\ref{app:empirical}; precise deployment
identifiers are omitted for anonymity.}---and
they are \emph{combinatorial}:
a solver executing several intents together can net opposing flows and
save on fees, producing more value than executing them separately; this
batching logic is central in MEV-aware market design
too~\cite{daian2020,canidio2023}. What
makes them unlike the combinatorial auctions of the literature is the
absence of a \emph{numeraire}. Users may demand any asset for any other,
there is no common unit of account, and assets are illiquid, so the
extra value created by batching cannot simply be paid out and shared.

A deployed intent-auction protocol uses the \emph{fair combinatorial
auction}: it constructs a fairness benchmark \emph{from the
bids themselves}: each trader's standalone competitive outcome defines a
floor, and a batched bid is executed only if it delivers at least that
floor to \emph{every} covered trader. This sidesteps the impossibility
of sharing batching gains by never allowing a batch to make any trader
worse off than individual competition would. The equilibrium behavior of
this design was analyzed by Canidio and Henneke~\cite{canidio2024}, who
identify a tension between the fairness guarantee and the value returned
to traders. Thus the model and benchmark are theirs; our contribution is
the worst-case approximation layer around them: how much
welfare does the fairness filter cost, when do competitive equilibria
exist, what does strategyproofness cost, and how do these prices behave
online and under manipulation.

\paragraph{A common scale.}
We organize such a theory around a recurring structural quantity. For a
solver with batch value $V_i(S)$ on a set $S$ of intents and standalone
values $V_i(\{j\})$, the \emph{complementarity factor} is the least
$g\ge1$ with $V_i(S)\le g\sum_{j\in S}V_i(\{j\})$ for all $i,S$: how much
more a batch can be worth than the sum of its parts. The benchmark floor
of order $j$ is $c_j=\max_i V_i(\{j\})$, and liquidity is
$\theta=1-\lambda\in[0,1]$, where $\lambda$ is the haircut on moving
value across assets ($\theta=0$ is the pure no-numeraire regime,
$\theta=1$ restores a numeraire). We show that $g$---together with
$\theta$---is the natural scale in a wide range of worst-case quantities,
which gives the design a unified analysis rather than a collection of
unrelated bounds.

The role of $g$ is deliberately worst-case: the theorems do not require a
deployment to estimate the structural complementarity of every batch.
The empirical section only shows that the corresponding observable
spreads and supports are real in deployment.

\paragraph{Contributions.}
We ask what price several design requirements carry, each in terms of $g$
(and liquidity $\theta$). Each item below leads with tight
characterizations (exact or asymptotically tight); partial frontiers and
peripheral extensions are marked as such.
\begin{enumerate}[leftmargin=*,itemsep=2pt]
\item \textbf{Price of fairness (Section~\ref{sec:pof}).} The welfare
ratio between the unconstrained optimum and best fair outcome is exactly
$\min(g,1/\theta)$, tight and order/price-independent, with endpoints
$\{g,1\}$; under coverage floor $\alpha$ the exact law is
$\min(g,\alpha+(1-\alpha)/\theta)$. The binding instance is an asymmetry
at the repair/drop indifference point, not maximal complementarity.

\item \textbf{Equilibrium existence and the core
(Section~\ref{sec:walras}).} A Walrasian/UDP equilibrium supporting an
efficient allocation exists for every rate-realizable floor-vector
instance iff $g=1$, with obligations equal to the fairness floors. The
deployed filter is the singleton-core constraint, the individual allocation
lies in the tight $g$-core, and the no-numeraire core can be nonempty when
the transferable core is empty.

\item \textbf{Price of strategyproofness (Section~\ref{sec:incentives}).}
Endogenous best-bid fairness conflicts with DSIC already on one order,
and online realized floors force ratio $g$ even randomized. In the public
historical-floor regime, strategyproofness costs $\Theta(1+\log g)$ for
many single-minded solvers and for one multi-minded solver. As further
separations, independent geometric option pricing admits an
$\Omega(g/\log^2 g)$ loss, and the multi-agent multi-minded cell remains a
partial frontier: it has an $O(N\log g)$ mechanism and a serial
$\Omega(\min(N,g))$ barrier.

\item \textbf{Online and strategic prices (Section~\ref{sec:online}).}
Online fair settlement pays exactly $g$ against realized floors; with
static public floors the residual one-way search problem is deterministic
$\Theta(\sqrt g)$ and randomized $\Theta(\log g)$. The individual
first-price restriction is smooth but leaves a latent-batch obstruction
for the full game; Appendix~\ref{app:manipulation} records a complete but
peripheral benchmark-manipulation extension with feedback
$g\mapsto g/(1-\mu)$.
\end{enumerate}

Fair winner determination is NP-hard but FPT in directed pairs.
Section~\ref{sec:empirics} reports diagnostics; the worst-case theory is
independent of them. Table~\ref{tab:summary} collects the results.

\begin{table}[t]
\centering
\caption{Worst-case prices as functions of the complementarity factor
$g$ and liquidity $\theta$ ($N$: number of strategic solvers; $\alpha$:
coverage). ``exact'' marks literal equalities or fully tight
characterizations; ``asymptotically tight'' marks $\Theta(\cdot)$
bounds.}
\label{tab:summary}
\footnotesize
\setlength{\tabcolsep}{2pt}
\begin{tabular}{@{}>{\raggedright\arraybackslash}p{0.27\textwidth}
                  >{\raggedright\arraybackslash}p{0.48\textwidth}
                  >{\raggedright\arraybackslash}p{0.20\textwidth}@{}}
\toprule
Quantity & Result & Status\\
\midrule
Price of fairness & $\min(g,1/\theta)$; endpoints $\{g,1\}$ & exact\\
\quad coverage refinement & $\min(g,\alpha+(1-\alpha)/\theta)$; size-irrelevant & exact\\
Walrasian/UDP existence & rate-realizable class: iff $g=1$ & exact\\
Core & filter $=$ singleton-core; $g$-core nonempty, tight & exact\\
Winner determination & NP-hard; FPT by directed-pair count & complexity classification\\
Strategyproofness (screening) & rand.\ $\Theta(1+\log g)$; det.\ exact $\sqrt g$ & tight/asympt.\ tight\\
\quad single-/multi-minded & both $\Theta(1+\log g)$; indep.\ geometric: $\Omega(g/\log^2 g)$ loss & asympt.\ tight / separation\\
\quad multi-agent multi-minded & $O(N\log g)$ / serial $\Omega(\min(N,g))$ & partial (serial barrier)\\
Online (realized floors) & exactly $g$, even randomized & exact\\
Online (static floors) & det.\ $\Theta(\sqrt g)$, rand.\ $\Theta(\log g)$ & asympt.\ tight\\
Price of anarchy & individual-only $\le e/(e-1)$; full conj.\ $\Theta(g)$ & partial (latent batches)\\
Contamination & threshold $\kappa_1\ge v\beta$ deters; $g\mapsto g/(1-\mu)$ & App.~\ref{app:manipulation}\\
\bottomrule
\end{tabular}
\end{table}

\paragraph{Related work.}
The model and the equilibrium analysis of the fair combinatorial auction
are due to Canidio and Henneke~\cite{canidio2024}; we complement their
$2\times2$ equilibrium study with a worst-case, approximation-theoretic
analysis. Recent intent-market and solver-competition models give
nearby strategic context~\cite{chitra2024}, and batch trading has been
proposed as a response to AMM arbitrage, LVR, and sandwich
attacks~\cite{canidio2023}. Combinatorial auctions are
classical~\cite{cramton2006}, but almost always with a numeraire and
quasilinear utilities; our setting has neither, putting it closer to
mechanism design without money~\cite{procaccia2013} and no-transfer
impossibility phenomena in MEV-aware fee mechanisms~\cite{bahrani2024}.
The price of fairness was introduced for divisible
resources~\cite{bertsimas2011} and developed for indivisible goods under
proportionality, envy-freeness, and
MMS~\cite{caragiannis2012,bei2019,barman2020}; here it scales with
production complementarity and liquidity, not agent count. Equilibrium
existence draws on assignment markets~\cite{shapley1971} and package
assignment~\cite{bikhchandani2002}. Our smoothness bound uses
Syrgkanis--Tardos~\cite{syrgkanis2013} and intrinsic
robustness~\cite{roughgarden2015}; online search instantiates one-way
trading~\cite{elyaniv2001}. Order-flow-auction models under PBS
provide related strategic context~\cite{ma2025}.

\section{Model}\label{sec:model}

There are $n$ trade intents (\emph{orders}) and a set of \emph{solvers}.
A solver $i$ submitting a batched bid on a set $S$ specifies, for each
$j\in S$, a \emph{raw output} $y_{i,S,j}\ge0$ in $j$'s requested asset
(trader-normalized units), before any within-batch rebalancing; the batch
value is
$V_i(S)=\sum_{j\in S}y_{i,S,j}$, realized only if the batch executes.
Delivery is batch-specific: the same solver may have different raw outputs
for order $j$ in different batches, so $V_i(S)$ is not additive in
general. The standalone value is $V_i(\{j\})=y_{i,\{j\},j}$, and the
floor is $c_j=\max_iV_i(\{j\})$. Reports are \emph{binding}: a solver
cannot claim raw output or delivered value it cannot produce, so relevant
misreports are downward. Welfare $W=\sum_jw_j$ is an additive social
objective over trader-normalized delivered scores, not a transferable
monetary surplus.

\begin{definition}[Complementarity factor]\label{def:g}
The instance has complementarity factor $g\ge1$ if
$V_i(S)\le g\sum_{j\in S}V_i(\{j\})$ for every solver $i$ and every $S$.
Thus $g$ caps how much batching amplifies standalone value; $g=1$ is
additive.
\end{definition}

\begin{definition}[Fairness floor; fair allocation]\label{def:fair}
The \emph{floor} of order $j$ is $c_j:=\max_i V_i(\{j\})$, the value of
its best standalone bid; $C:=\sum_j c_j$. An allocation is \emph{fair} if
every allocated order receives $w_j\ge c_j$; batched bids violating this
for some covered order are \emph{filtered}.
\end{definition}

An allocation selects winning bids covering disjoint order sets; its
welfare is $W=\sum_j w_j$. Let $\OPT$ be the maximum welfare ignoring
fairness and $\FAIR$ the maximum over fair allocations; the \emph{price
of fairness} is $\PoF:=\sup\OPT/\FAIR$ over a stated class.

\begin{definition}[Liquidity and feasible delivery]\label{def:theta}
Within a single winning bid a solver may convert value across its traders
at haircut $\lambda\in[0,1]$. Write $\theta:=1-\lambda$ for
\emph{liquidity}: moving one unit out of one trader's asset yields
$\theta$ units in another. For delivered vector $w$, let
$D^+=\sum_{w_j>y_{i,S,j}}(w_j-y_{i,S,j})$ and
$D^-=\sum_{w_j<y_{i,S,j}}(y_{i,S,j}-w_j)$; delivery is feasible iff
$D^+\le\theta D^-$.
Thus a starved leg may receive $w_j>y_{i,S,j}$, funded by over-served legs;
$\theta=1$ is value-preserving and $\theta=0$ forbids cross-trader repair.
Conversion is per-bid.
\end{definition}

This per-bid model excludes cross-settlement inventory; richer liquidity
only eases repair, so $\min(g,1/\theta)$ is a worst-case ceiling.

The deployed mechanism additionally requires uniform clearing prices per
directed token pair (\emph{uniform directed prices}, UDP), forcing all
orders on one directed pair into a single winning bid. UDP is vacuous for
the constructions in Section~\ref{sec:pof}, which use distinct pairs.
Sections~\ref{sec:incentives}--\ref{sec:online} study public-benchmark
screening and online abstractions without imposing UDP. Section~\ref{sec:walras}
uses UDP as its obligation abstraction, and UDP later reappears as the
source of fixed-parameter tractability. This is the order-level analogue of the uniform clearing
price in frequent batch auctions~\cite{budish2015}, specialized to the
trade-intent model of Canidio and Henneke~\cite{canidio2024}. Full
proofs of all numbered results are in the appendix.

\paragraph{Regime map.}
\vspace{-0.5em}
\begin{center}
\scriptsize
\renewcommand{\arraystretch}{0.82}
\setlength{\tabcolsep}{2pt}
\begin{tabular}{@{}p{0.22\textwidth}p{0.23\textwidth}p{0.12\textwidth}p{0.07\textwidth}p{0.25\textwidth}@{}}
\toprule
Result & Floors & Repair & UDP & Solvers\\
\midrule
Thms.~\ref{thm:pof}--\ref{thm:coverage} & standalone/given & yes, $\theta$ & no & any\\
Thm.~\ref{prop:walras} & endogenous & no & yes & many\\
Thm.~\ref{thm:incompat} & current-report & no & no & many\\
Thms.~\ref{thm:posp}--\ref{thm:lift} & public/historical & no & no & screening; 1 multi; many single\\
Thm.~\ref{thm:multi} & public/historical & no & no & $N$ multi\\
Prop.~\ref{prop:online}--Thm.~\ref{thm:online} & realized/static & no & no & online\\
Thm.~\ref{thm:poa} & given & no & no & many\\
\bottomrule
\end{tabular}
\end{center}
\vspace{-0.7em}
Except in §3, repair is not part of feasible deviations.

\section{The Price of Fairness}\label{sec:pof}

The first price is informational: once the benchmark floors are fixed by
standalone competition, how much welfare can the componentwise fairness
filter destroy?  The answer is governed exactly by $g$ and liquidity
$\theta$.

\begin{theorem}[Price of fairness, exact]\label{thm:pof}
For every $\theta\in[0,1]$ and $g\ge1$, over all instances with liquidity
$\theta$ and complementarity factor at most $g$,
\[
\PoF(\theta,g)=\min\!\big(g,\,1/\theta\big),
\]
tight, and independent of the number of orders and of prices
($1/\theta:=\infty$ at $\theta=0$).
\end{theorem}

\begin{proof}[Proof idea]
\emph{Upper bound.} Process each bid of an unconstrained optimum
independently. For a locally undominated OPT bid, replacing it by
singleton floor allocations would not improve welfare, so
$\gamma=V/C_S\ge1$; Definition~\ref{def:g} gives $\gamma\le g$. Its
deficit is $D=\sum_{j\in S}(c_j-y_{i,S,j})^+\le C_S$. Either drop the batch and settle
floors, delivering $C_S=V/\gamma$, or repair starved legs by a feasible
delivery that saturates the conversion budget of
Definition~\ref{def:theta}. Clearing $D$ loses $D(1-\theta)/\theta$, so
whenever repair is feasible,
\[
  \frac{V-D(1-\theta)/\theta}{V}
  \ge 1-\frac{1-\theta}{\theta\gamma}
  \ge \theta .
\]
Dropping covers low gains by $V/g$; repair covers high gains by
$V/(1/\theta)$; summing over disjoint optimal bids gives the bound.
\emph{Lower bound.} A two-order instance with one starved leg, calibrated
to $\gamma^\ast=\min(g,1/\theta)$, forces
$\FAIR=c$ against $\OPT=\gamma^\ast c$. The binding instance is an
extreme asymmetry, not maximal complementarity.\hfill$\square$
\end{proof}

\begin{corollary}[Dichotomy]\label{cor:dichotomy}
At the endpoints, the pure no-numeraire auction pays the full factor
$g$, while a full numeraire removes the fairness cost entirely.
\end{corollary}

A natural attempt to add a third parameter fails, and identifies the
right one.

\begin{proposition}[Batch size is not the parameter]\label{prop:dirrel}
With $\PoF(\theta,g,d)$ the price restricted to bids covering at most
$d$ orders,
\[
  \PoF(\theta,g,1)=1,\qquad
  \PoF(\theta,g,d)=\min(g,1/\theta)\quad(d\ge2).
\]
\end{proposition}

Say an instance has \emph{coverage} $\alpha\in[0,1]$ if every bid
satisfies $y_{i,S,j}\ge\alpha c_j$ on every covered order---no leg is
worse than $\alpha$ times the benchmark, a per-leg routing-quality floor.

\begin{theorem}[The coverage law]\label{thm:coverage}
Write
$\Gamma(\theta,\alpha):=\alpha+(1-\alpha)/\theta$, with
$\Gamma(0,\alpha)=\infty$ for $\alpha<1$ and $\Gamma(0,1)=1$. Over
instances with liquidity $\theta$, complementarity at most $g$, and
coverage $\alpha$,
\[
\PoF(\theta,g,\alpha)=\min\!\big(g,\Gamma(\theta,\alpha)\big),
\]
tight on two-order instances.
\end{theorem}

The quantity $\Gamma-1=(1-\alpha)\cdot\frac{1-\theta}{\theta}$ is the
largest relative per-leg shortfall times the conversion loss per unit of
delivered deficit: the price of fairness prices the worst starved leg,
discounted by liquidity. The endpoint $\alpha=1$ is exactly the
componentwise filter-passing under which fairness is free. This
refinement is also why batch size disappears from the theorem: what
matters is not how many orders are bundled, but how much one covered leg
can fall below its own standalone benchmark before the batch's surplus
can be converted back to repair it.

\section{Equilibrium Existence and the Core}\label{sec:walras}

The second question is structural. When does the same filter coincide
with a competitive equilibrium or a cooperative stability notion rather
than merely a feasibility rule? Here $g=1$ is the threshold.
This is not a restatement of the gross-substitutes existence theory:
gross substitutes support Walrasian equilibria with prices and transfers
over items~\cite{kelso1982,gul1999}, whereas our frontier is about
componentwise obligations in traders' requested assets.
Throughout this section, equilibrium feasibility is the native UDP
obligation model: an executable deviation must meet each obligation
componentwise, equivalently $\theta=0$ for deviations. The liquidity repair
operator of Definition~\ref{def:theta} is used for the price-of-fairness
comparison, not for defining executable UDP deviations.

We work in the order-price abstraction of UDP: a directed-rate system
induces order-level obligations $w_j\ge0$ that any solver serving $j$
must deliver, and the single property used is that $w_j$ depends only on
the order. A \emph{UDP equilibrium supporting the efficient allocation}
is a rate system and an allocation in which every order is served at
exactly its $w_j$, every solver's executed portfolio maximizes its
total profit
$\sum_{S\in\mathcal P}\sum_{j\in S}(y_{i,S,j}-w_j)$ over feasible
portfolios $\mathcal P$ (those with $y_{i,S,j}\ge w_j$ for every
executed bid $S$ and covered order $j$), and produced value equals
$\OPT$.
This definition deliberately separates prices from payments: $w_j$ is an
obligation in order $j$'s normalized delivered units, not a transfer that
can compensate another trader. That separation is what makes the
frontier different from the transferable package-assignment theorem.

\begin{lemma}[Batch domination under $g=1$]\label{lem:domination}
If $g=1$ then at any rates every executable batched bid is weakly
dominated by the sub-portfolio of its positive-profit singletons:
$\sum_{j\in S}(y_{i,S,j}-w_j)=V_i(S)-\sum_jw_j\le
\sum_j(V_i(\{j\})-w_j)^+$.
Hence every solver has an individual-only best response; the one use of
UDP is that $w_j$ is the same in a batch or a singleton.
\end{lemma}

The lemma is the reason $g=1$ is special. If batching is never more
productive than the sum of singletons, then any profitable executable
batch contains enough profitable singleton mass to replace it. The
no-numeraire feasibility constraint only strengthens this conclusion:
some batch deviations that would be available with transfers are simply
not executable componentwise.
Conversely, once $g>1$, overlap matters. A solver can create value on a
pair without creating enough value on either order alone to support
order-level obligations. The three-cycle lower bound exploits exactly
this gap: each pair wants one obligation low for feasibility and another
high for incentive compatibility, and the cycle leaves no price vector
that satisfies all three demands.

\begin{theorem}[Existence frontier is $g=1$]\label{prop:walras}
Within the order-price abstraction of UDP, restricted to instances whose
floor vector $c$ is rate-realizable, a UDP equilibrium supporting an
efficient allocation exists for every instance with complementarity at
most $g$ iff $g=1$. For $g=1$, the obligations $w_j=c_j$ form a
best-bid-fair equilibrium. For every $g>1$, some rate-realizable instance
admits no such equilibrium. Solvers have no capacity constraints.
\end{theorem}

\begin{proof}[Proof idea]
\emph{$g=1$:} set $w_j=c_j$ and assign each order to a best producer;
winners earn zero margin, losers nonpositive, and
Lemma~\ref{lem:domination} kills batch deviations. Produced value is
$C=\OPT$, and the price vector coincides with the fairness benchmark.
\emph{$g>1$:} the obstruction is native---non-existence with transfers
does not imply it here, since feasibility only \emph{removes}
deviations. The hard instance is a three-cycle of overlapping
pair-batches with vector $(1+\tfrac\delta2,1+\tfrac\delta2)$ over floors
$1$. The three orders lie on distinct directed pairs, so the unit floor
vector is rate-realizable. Supporting one efficient pair forces an
obligation on the third order low enough to keep that pair feasible. The
adjacent unchosen pair is
then both executable and must be unprofitable; splitting on whether the
shared obligation is above or below its singleton floor yields either
$\delta\le0$ or a componentwise-feasibility contradiction. Full details are
in Appendix B.\hfill$\square$
\end{proof}

The filter also has a cooperative reading. A bid \emph{blocks} an
allocation at factor $\kappa$ if some feasible delivery gives every
covered order strictly more than $\kappa w_j$; the $\kappa$-core forbids
this.

\begin{proposition}[Core placement]\label{prop:core}
(i) Best-bid fairness is exactly immunity to \emph{singleton} blocking at
factor $1$: the deployed filter is the singleton-coalition core
constraint. (ii) The individual allocation $w_j=c_j$ lies in the
$g$-core for every $\theta$, and $g$ is the least universal
blocking-immunity factor. (iii) There are instances whose
transferable-utility core is empty while the $\theta=0$ core is
nonempty---a separation driven by overlap, like the integrality gap.
\end{proposition}

Thus the fairness filter is not an ad hoc admissibility test; it is the
smallest coalition-stability condition visible when coalitions are
restricted to orders' own requested assets. The multiplicative $g$-core
statement says that even when exact support fails, the individual
allocation is never blocked by more than the same complementarity factor
that drives Section~\ref{sec:pof}. This is one of three appearances of a single
phenomenon: the no-numeraire constraint \emph{removes deviations} (VCG
payments become infeasible; batch deviations are dominated at $g=1$;
componentwise blocking is weak). Whether the $\theta=0$ core is always
nonempty is open: Scarf's route needs balancedness, broken by the same
overlap geometry.

\paragraph{Complexity.}
Fair winner determination is NP-hard (from 3-dimensional matching) but
fixed-parameter tractable in the number $P$ of directed asset pairs:
Under UDP, all orders on the same directed pair settle atomically: a
feasible solution either assigns the entire directed-pair block to one
bid or leaves it unserved. Thus every fair bid is a weighted subset of
the $P$ atoms. Dynamic programming over pair masks
(or equivalently weighted set packing on $P$ items) tries
$2^P$ states and scans the bid list in each state, giving
$2^{O(P)}\cdot\mathrm{poly}(|\mathcal B|)$.
A deployment cap on pairs per batch thus makes the problem tractable in
the intended regime. This is another way in which $g$ and the
token-pair graph play complementary roles: $g$ controls the loss from
fairness and incentives, while $P$ controls the search space left by the
uniform-price constraint.

\section{The Price of Strategyproofness}\label{sec:incentives}

The third price is strategic. Truthfulness is not studied under the
original endogenous benchmark, because the preceding structure creates a
trilemma: current-report floors, ex-post fairness, and VCG-style
strategyproofness cannot coexist without transfers. This places the
problem in the no-money mechanism-design tradition~\cite{procaccia2013},
with a blockchain-specific impossibility flavor close to post-MEV fee
mechanism design~\cite{bahrani2024}.

Throughout this section, incentive compatibility means \emph{one-sided
dominant-strategy implementability under binding reports}: upward
misreports are infeasible by settlement, and ``universal truthfulness''
refers to posted-demand mechanisms whose menu is independent of the
solver's behavior.

There are three benchmark regimes. First, \emph{current-report} floors are
the original endogenous ideal, but they are self-referential:
VCG-style compensation is infeasible because the no-numeraire constraint
gives no asset in which to make the Clarke payment. More sharply,
best-bid fairness and dominant-strategy truthfulness already conflict on
one order:

\begin{theorem}[Best-bid fairness is incompatible with strategyproofness]
\label{thm:incompat}
Call a mechanism \emph{non-wasteful} if it allocates an order whenever
some reported standalone value is positive, and \emph{value-respecting}
if it allocates to a solver with maximal reported value. No feasible,
non-wasteful, value-respecting, DSIC mechanism guarantees fairness at the
endogenous best-bid benchmark $c_j=\max_i r_{ij}$. This holds already for
one order. Under DSIC, the enforceable benchmark collapses to the
runner-up level.
\end{theorem}

Second, \emph{realized lifetime} floors are operationally stronger, but
as Section~\ref{sec:online} proves, online fairness then admits no ratio
better than $g$ even with randomization: a solver settling a batch can be
undercut by a later standalone bid that raises a covered order's own
floor. Third, \emph{public delayed or historical} floors sever this
self-reference while preserving the spread $g$ as the operative parameter.
This is the deployed-safe regime, not a hypothetical substitute: a
deployment-safe benchmark is already a public historical window. We therefore
analyze this exogenous-benchmark regime; the price of strategyproofness
$\mathrm{PoSP}$ is the welfare ratio between $\FAIR$ and the best
universally truthful fair mechanism, and solvers are single- or
multi-minded with private values passing the filter. Under endogenous
floors the relevant question is not the price of strategyproofness but its
non-existence.

The canonical screening family of the regime is already nontrivial for
one strategic solver:

\begin{theorem}[Canonical screening family]\label{thm:posp}
Over the screening family $R_g$ (one strategic solver, public floors),
the price of strategyproofness is $\Theta(1+\log g)$: every randomized
DSIC mechanism pays at least $\frac{g\ln g}{2(g-1)}\ge\frac{\ln g}{2}$,
a geometric posted-price mechanism achieves $e(\lceil\ln g\rceil+1)$,
and the best deterministic mechanism is exactly $\sqrt g$.
\end{theorem}

The lower bound is a model-native technique: a one-sided envelope
inequality (downward incentive compatibility alone forces monotone
utilities, so the utility derivative dominates the execution probability
a.e.) integrated against a $V^{-2}$ kernel. In the screening family a
type $V$ can only underreport; if $p(V)$ is execution probability and
$u(V)$ interim utility, downward IC gives
\[
  u(V)\ge u(V')+(V-V')p(V') \qquad (V'\le V),
\]
so $u'(V)\ge p(V)$ a.e.; the scale-free integral forces logarithmic loss.
The deterministic optimum balances the two errors at threshold $\sqrt g$.
Random geometric posted levels give the matching upper bound up to
constants. Two mechanisms lift the guarantee to the full game.

\paragraph{Mechanism model.}
In the exogenous-benchmark regime the floors $c_j$ are public. The public
geometry consists of feasible option sets $S$ and their floor masses
$C_S=\sum_{j\in S}c_j$. A solver's private type specifies, for each
option $S$, an executable filter-passing production vector
$y_{S,j}\ge c_j$; we summarize it by the scalar value
$V_S=\sum_{j\in S}y_{S,j}\in[C_S,gC_S]$. Any total delivery
$D\in[C_S,V_S]$ can clear all floors componentwise by scaling surplus:
deliver $w_j=c_j+\alpha(y_{S,j}-c_j)$ for
$\alpha=(D-C_S)/(V_S-C_S)$ (with the obvious convention when
$V_S=C_S$). The mechanisms below post only as a function of public
$C_S$ and a public random seed, and each solver responds by demanding a
disjoint collection of options rather than reporting values. Settlement
verifies the promised componentwise delivery. Universal truthfulness is
in this posted-demand sense: for every seed, demanding a
utility-maximizing executable collection is dominant because the menu is
independent of that solver's behavior.
The posted quantity $C_St$ is not a monetary payment; it is a required
delivered-score threshold for option $S$, in the covered orders' requested
assets. The profit $V_S-C_St$ is retained raw value, not cash.

\begin{theorem}[Two routes to $\Theta(1+\log g)$]\label{thm:lift}
(i) For arbitrarily many single-minded solvers with public floors, an
eligibility-posted packing mechanism is universally truthful and
$\Theta(1+\log g)$-competitive. (ii) For one multi-minded solver of
arbitrary structure, a coupled-menu mechanism---post all options at
prices $C_S\,t$ for a single random multiplier $t=e^k$ and let the solver
demand any profit-maximizing disjoint collection---is universally
truthful and $\Theta(1+\log g)$-competitive. The natural independent
geometric per-option analogue can instead pay $\Omega(g/\log^2 g)$.
\end{theorem}

The multi-minded bound rests on a \emph{profit-curve} technique: the
solver's profit $p(t)$ over disjoint collections is convex and piecewise
linear, delivered mass is $t\,(-p'(t))$, and a geometric grid telescopes
$p(1)$ against expected welfare. The single random multiplier is
essential: the independent geometric analogue lets a multi-minded solver
take only the cheapest member of a family of substitutes, losing the
collection value that fairness must compare against. The remaining cell
is both multi-agent and multi-minded:

\begin{theorem}[Multi-agent multi-minded]\label{thm:multi}
Random-priority coupled menus are universally truthful with
$\E[W]\ge\FAIR/\big(N(e-1)(\lceil\ln g\rceil+1)+1\big)$, i.e.
$\Theta_N(1+\log g)$ for a constant number $N$ of strategic solvers.
Conversely every \emph{serial-demand} mechanism---a report-independent
ex-ante service order, arbitrary per-identity multiplier distributions,
each agent demanding a profit-maximizing collection---pays
$\Omega(\min(N,g))$, even with geometry-aware ordering. So any
$N$-independent $\mathrm{polylog}(g)$ mechanism must break the
serial-demand format.
\end{theorem}

The barrier is proved by a \emph{symmetrized swarm}: all identities share
the same public geometry (all singletons plus the grand set), differing
only in private values, so no report-independent order can tell the
batcher from the blockers. The telescoping identity
\[
  \sum_i(1-q_i)\prod_{s\prec i}q_s=1-\prod_s q_s\le1
\]
then places the batcher behind an expected blocker.

These results should be read as a trichotomy. Single-minded competition
is handled by posted eligibility and packing; one multi-minded solver is
handled by a coupled menu and profit-curve telescoping; many
multi-minded solvers expose a genuine serial-demand barrier. The open
question is whether a mechanism that is not serial-demand can remove the
dependence on $N$ while keeping the logarithmic dependence on $g$.

\section{Online and Robustness}\label{sec:online}

The remaining results explain why deployed systems replace current-report
floors by historical, public, and rate-limited ones.

\paragraph{Online floors.}
With realized lifetime floors, online fairness can be attacked by a later
standalone bid that raises a covered order's own floor. This is the
online analogue of the endogenous-benchmark impossibility in
Section~\ref{sec:incentives}: the solver competes against a future
revision of the fairness constraint itself.
\begin{proposition}[Realized floors force a factor $g$]\label{prop:online}
The deterministic online competitive ratio of fair settlement is exactly
$g$.
\end{proposition}
\begin{theorem}[Two online prices]\label{thm:online}
(i) Against realized lifetime floors with almost-sure fairness,
randomization does not help: the competitive ratio is exactly $g$. (ii)
With static public floors on single-contention instances, the residual
problem is one-way search over the spread, with deterministic ratio
$\Theta(\sqrt g)$ and randomized ratio $\Theta(\log g)$; ratio-threshold
greedy mechanisms still pay $\Omega(g)$.
\end{theorem}
Thus Theorem~\ref{thm:online} is the operational reason for the
exogenous-benchmark regime in Section~\ref{sec:incentives}; once floors
are fixed, the residual problem is one-way search over width
$g$~\cite{elyaniv2001}.

\paragraph{Strategic play.}
For the first-price game, measured in produced welfare $\PW$ rather than
delivered welfare, the individual-only restriction is smooth in the
standard composable-mechanisms sense~\cite{syrgkanis2013,roughgarden2015}:
\begin{theorem}[Individual competition is smooth]\label{thm:poa}
Every coarse correlated equilibrium of the per-order-decomposable
individual-only game has
$\E[\PW]\ge(1-\tfrac1e)\OPT_{\mathrm{ind}}$, hence price of anarchy at
most $e/(e-1)$ relative to $\OPT_{\mathrm{ind}}$ and at most
$\frac{e}{e-1}\,g$ relative to the unrestricted optimum. The full-game
price of anarchy is conjectured $\Theta(g)$; the obstruction is latent
batches that execute only under a deviation.
\end{theorem}

\paragraph{Benchmark robustness.}
Benchmark manipulation is itself a priced game: a sufficient injection-cost
floor deters it, and steady contamination feeds back as
$g\mapsto g/(1-\mu)$, closing the loop with every floor-calibrated bound
(Appendix~\ref{app:manipulation}).

\section{Deployment Diagnostics (Motivation)}\label{sec:empirics}

An anonymized trace from a deployed intent-auction system contains 802,816 post-adoption solver
competitions over eight months. On the sanity-filtered USD subsample it
reports \$54.69B covered notional and a \$5.13M visible accounting gap
(0.94 bps). The visible spread is nontrivial in the tail:
for all auctions, $\widehat G^{\mathrm{obs}}>2$ in 25.3\% and
$\widehat G^{\mathrm{obs}}>4$ in 20.2\%; for small-notional auctions,
$\widehat G^{\mathrm{obs}}>2$ in 63.1\%. The
appendix defines the score variables, filters, and replay diagnostics.
When the denominator is positive, the parser computes
$\widehat G^{\mathrm{obs}}_a=Score_{1,a}/Score^{\mathrm{other}}_{2,a}$,
where $Score^{\mathrm{other}}_{2,a}$ is the best positive valid score
after removing the winning solver, not a singleton floor mass. Thus the
spread is not an estimate of $g$, but it gives a conditional lower
envelope. Let $S_a$ be the winning batch and
$C_a=\sum_{j\in S_a}c_j$ its floor mass. If
$Score^{\mathrm{other}}_{2,a}=\eta_a C_a$, then
$\eta_a\widehat G^{\mathrm{obs}}_a=Score_{1,a}/C_a\le g_a$, because
$Score_{1,a}\le V_{\mathrm{win}}(S_a)\le
g_a\sum_{j\in S_a}V_{\mathrm{win}}(\{j\})\le g_aC_a$.
The trace does not observe $C_a$, so $\eta_a$ is an auxiliary-audit
ratio; without certifying $\eta_a\ge1$ we do not read
$\widehat G^{\mathrm{obs}}$ as an unconditional lower bound on $g$.
These diagnostics motivate the parameter regime; they neither estimate
structural $g$ nor test point predictions. The bounds are tight as
worst-case functions of $g$ and hold for whatever $g$ a deployment
exhibits. The trace also records $P$ and a reserve diagnostic: 97.1\% of
auction-weighted cell-windows are vulnerable under a naive local hard
reserve, motivating inherited, quantile-based, delayed, and rate-limited
benchmarks.

\section{Conclusion}\label{sec:conclusion}

A recurring scale---the complementarity factor $g$, with liquidity
$\theta$---organizes the worst-case fairness, equilibrium, complexity,
incentive, and online prices of the no-numeraire fair combinatorial
auction. Two phenomena recur: no-numeraire feasibility
removes some deviations while disabling transfers, and one geometric
family prices both informational and incentive costs of the spread.

Three questions remain open, each with an identified obstruction:
full-game price of anarchy (conjectured $\Theta(g)$; the latent-batch
obstruction), polylogarithmic mechanisms for multi-agent multi-minded
solvers without $N$ dependence (the serial-demand barrier), and nonemptiness of the
$\theta=0$ core (Scarf balancedness, broken by overlap).

\bibliographystyle{splncs04}
\bibliography{refs}

\appendix
\section{Full Proofs for Section~\ref{sec:pof}}

\begin{proof}[of Theorem~\ref{thm:pof}]
\emph{Endpoint.} If $\theta=0$, repair is unavailable; dropping every OPT
batch to the singleton floor allocation loses at most the factor $g$, and
the lower-bound construction below attains $g$. Hence assume $\theta>0$
for the repair calculation.

\emph{Upper bound.} Fix an unconstrained optimum and process its winning
bids independently; bids cover disjoint order sets, so the repaired
welfare contributions add. Take a winning bid with floor mass $C_S=\sum_{j\in S}c_j>0$
(a bid with $C_S=0$ has all floors zero and is already fair), value
$V=\sum_{j\in S}y_j$, where $y_j:=y_{i,S,j}$, gain
$\gamma=V/C_S$, and deficit $D=\sum_{j\in S}(c_j-y_j)^+$. By local
optimality of $\OPT$, we may assume $V\ge C_S$ (otherwise replacing this
bid by the singleton floor allocation on $S$ improves welfare), so
$\gamma\ge1$; Definition~\ref{def:g} gives $\gamma\le g$. Two fair uses:

\emph{Repair.} Move value from over-served legs to starved ones. The
total surplus available is $s=V-(C_S-D)=V-C_S+D$, and clearing the
deficit costs $D/\theta$ units of surplus (haircut $\lambda=1-\theta$).
This is feasible iff $V-C_S\ge D\frac{1-\theta}{\theta}$, and then the
delivery saturating the conversion budget of Definition~\ref{def:theta}
has welfare $f=V-D\frac{1-\theta}{\theta}$. Since every raw output
value is nonnegative, $(c_j-y_j)^+\le c_j$ and therefore $D\le C_S$.
Together with $V-C_S=(\gamma-1)C_S$, feasibility holds whenever
$\gamma-1\ge\frac{1-\theta}{\theta}$, i.e. $\gamma\ge1/\theta$; then
\[
  \frac{f}{V}
  =1-\frac{1-\theta}{\theta\gamma}\frac{D}{C_S}
  \ge 1-\frac{1-\theta}{\theta\gamma}
  \ge \theta
  = \frac{1}{1/\theta}.
\]
Thus repair gives $f\ge V/(1/\theta)$ whenever $\gamma\ge1/\theta$, and
hence $f\ge V/\!\min(g,1/\theta)$ in the repair regime.

\emph{Drop.} Settle every covered order at its floor: $f=C_S=V/\gamma$.
For $\gamma\le\min(g,1/\theta)$ this is $\ge V/\!\min(g,1/\theta)$.

Either way each bid retains a $1/\min(g,1/\theta)$ fraction of its value;
summing, $\FAIR\ge\OPT/\min(g,1/\theta)$.

\emph{Lower bound.} Let $\gamma^\ast=\min(g,1/\theta)$. Two orders on
distinct pairs, $c_1=0$, $c_2=c$; one batched bid with
$y_2=0$, $y_1=\gamma^\ast c$, so $V=\gamma^\ast c$, $C_S=c$, and the
batching solver's standalone values are $0,c$, so
Definition~\ref{def:g} holds as $\gamma^\ast\le g$. The unconstrained
optimum runs this batch, delivering $\gamma^\ast c$ (all on order~1), so
$\OPT=\gamma^\ast c$. Fair repair of order $2$ to its floor $c$ requires
withdrawing $c/\theta$ from $y_1$; if $\gamma^\ast=1/\theta$ this exhausts
$y_1$ exactly (delivered welfare $c$), and if
$\gamma^\ast=g<1/\theta$ the batch is irreparable (welfare $c$ from the
individual allocation). In both cases $\FAIR=c$ and
$\OPT=\gamma^\ast c$, so $\OPT/\FAIR=\gamma^\ast$. Independence of the
number of orders and of prices is immediate from the construction.\hfill$\square$
\end{proof}

\begin{proof}[of Proposition~\ref{prop:dirrel}]
The two-order family of Theorem~\ref{thm:pof} already has $d=2$ and
attains $\min(g,1/\theta)$, and the upper bound never uses $d$; for
$d=1$ only singletons exist, so $\OPT=\FAIR=C$.\hfill$\square$
\end{proof}

\begin{proof}[of Theorem~\ref{thm:coverage}]
\emph{Endpoint.} If $\theta=0$ and $\alpha<1$, repair is unavailable and
the drop argument gives ratio $g$; if $\alpha=1$, every admitted bid
already clears all floors and the ratio is $1$. Hence assume $\theta>0$
below.

\emph{Upper bound.} Write
$\Gamma=\Gamma(\theta,\alpha):=\alpha+(1-\alpha)/\theta$, with the
endpoint convention from the theorem. As above, coverage caps the deficit by
$D\le(1-\alpha)C_S$. Repair is feasible when
$\gamma-1\ge(1-\alpha)\frac{1-\theta}{\theta}=\Gamma-1$, and then
$f=V-\frac{D(1-\theta)}{\theta}\ge V\big(1-\frac{(1-\alpha)(1-\theta)}
{\theta\gamma}\big)\ge V/\Gamma$ for $\gamma\ge\Gamma$ (monotone in
$\gamma$, equality at $\gamma=\Gamma$); dropping gives $V/\gamma\ge
V/\Gamma$ for $\gamma\le\Gamma$ and $V/g$ for $\gamma\le g$. Hence
$\FAIR\ge\OPT/\min(g,\Gamma)$. \emph{Lower bound.} Two orders, $c_1=0$,
$c_2=c$; batched bid with $y_2=\alpha c$ (coverage saturated) and
$y_1=(\gamma^\ast-\alpha)c$, $\gamma^\ast=\min(g,\Gamma)$; the batching
solver's standalone values are $0,c$, so Definition~\ref{def:g} holds as
$\gamma^\ast\le g$. Repairing
order $2$ needs $(1-\alpha)c/\theta$ from $y_1$; if $\gamma^\ast=\Gamma$
then $y_1=(\Gamma-\alpha)c=\frac{(1-\alpha)c}{\theta}$ exactly, and if
$\gamma^\ast=g<\Gamma$ the batch is irreparable. Either way $\FAIR=c$,
$\OPT=\gamma^\ast c$.\hfill$\square$
\end{proof}

\section{Full Proofs for Section~\ref{sec:walras}}

\begin{proof}[of Lemma~\ref{lem:domination}]
\[
\begin{aligned}
\sum_{j\in S}(y_{i,S,j}-w_j)
&=V_i(S)-\sum_{j\in S}w_j\\
&\le \sum_{j\in S}V_i(\{j\})-\sum_{j\in S}w_j\\
&\le \sum_{j\in S}(V_i(\{j\})-w_j)^+ .
\end{aligned}
\]
The middle step uses $g=1$. Each
positive-profit singleton has $V_i(\{j\})>w_j$, hence is executable; a
non-executable batch is not a deviation. The obligation $w_j$ is the same
whether $j$ is served in the batch or individually, which is the only
property of UDP used.\hfill$\square$
\end{proof}

\begin{proof}[of Theorem~\ref{prop:walras}]
\emph{$(a)$, $g=1$.} Set $w_j=c_j$ and assign each order $j$ to a solver
attaining $V_i(\{j\})=c_j$. The winner earns zero margin per order; any
losing singleton has $V_i(\{j\})-c_j\le0$; and
Lemma~\ref{lem:domination} kills batch deviations (no positive-profit
singletons). With $g=1$, any allocation produces at most
$\sum_{(i,S)}V_i(S)\le\sum_j c_j=C$, attained here, so produced value is
$C=\OPT$, all delivered. The point $w_j=c_j$ is the best-bid floor, so
the equilibrium is fair (the trader-optimal core point of the underlying
assignment market~\cite{shapley1971}, to which
Lemma~\ref{lem:domination} reduces the economy).

\emph{$(b)$, $g>1$.} We argue natively; non-existence with transfers does
not carry over, as feasibility only removes deviations. Fix
$\delta\in(0,2(g-1)]$; orders $1,2,3$ on distinct pairs, so the floor
vector $c_j=1$ is rate-realizable. A rival has
standalone value $1$ on each ($c_j=1$), and three solvers each holding
one batched bid on a pair $\{j,j'\}$ with vector
$(1+\tfrac\delta2,1+\tfrac\delta2)$ (value $2+\delta\le2g$) and
standalone values $1$. The efficient allocation runs one batch, say
$\{1,2\}$, plus order $3$ individually (produced $3+\delta$). Suppose
rates $w$ support it. Order $3$ is served by a standalone solver of value
$1$, so $w_3\le1$; orders $1,2$ are served by the batch, so
$w_1,w_2\le1+\tfrac\delta2$. The batch $\{2,3\}$ is then executable, and
equilibrium requires it unprofitable against its solver's individual
alternatives: $(2+\delta)-w_2-w_3\le(1-w_2)^++(1-w_3)^+$. If $w_2\le1$
the right side is $(1-w_2)+(1-w_3)$, giving $\delta\le0$; if $w_2>1$ it
is $1-w_3$, giving $w_2\ge1+\delta>1+\tfrac\delta2$, contradicting
feasibility. No supporting rates exist.\hfill$\square$
\end{proof}

\begin{proof}[of Proposition~\ref{prop:core}]
(i) If $w_j<c_j$, the standalone bid attaining $c_j$ blocks $j$ at factor
$1$; if $w_j\ge c_j$ for all $j$, no singleton strictly improves. (ii)
Blocking $w_j=c_j$ at factor $g$ needs delivery $>g\,c_j$ to every
$j\in S$, totaling $>g\,C_S$, while any feasible delivery from a bid
totals $\le V_i(S)\le g\,C_S$ (haircuts only lose value; if all floors in
$S$ are zero then $V_i(S)=0$ and strict blocking is impossible).
Tightness: a batch with vector $(g'c,g'c)$ over floors $(c,c)$,
$g'\in(\kappa,g]$, blocks at any $\kappa<g$. (iii) Three-cycle floors
$(1,1,1)$, pair vectors $(1+\tfrac\delta2,1+\tfrac\delta2)$: summing the
three transferable core constraints $x_j+x_{j'}\ge2+\delta$ gives
$\sum x_j\ge3+\tfrac{3\delta}2>3+\delta=\OPT$, so the transferable core is
empty; the allocation (batch $\{1,2\}$, order $3$ individually) is
$\theta=0$-unblocked since every other batch caps a covered order at a
value already attained.\hfill$\square$
\end{proof}

\paragraph{Complexity.}
NP-hardness is a reduction from 3-dimensional matching: orders are ground
elements, each element has an individual bid of value $1$, and each triple
is represented by a solver whose singleton values on the three covered
elements are all $1$ and whose triple bid produces $1+\delta/3$ on each of
the three orders (total value $3+\delta\le3g$ for $\delta\le3(g-1)$). Thus
the triple satisfies the complementarity bound and passes the fairness
filter.
If the maximum number of disjoint triples is $m$, the optimum welfare is
$n+\delta m$, where $n$ is the number of ground elements; hence welfare
$n+\delta n/3$ is attainable iff a perfect matching exists. For
fixed-parameter tractability, treat each directed pair as an atomic item
under UDP: all orders on the same directed pair settle together, so a
feasible solution either assigns the entire directed-pair block to one
bid or leaves it unserved. After fair filtering, every surviving bid is a
weighted subset of the $P$ atoms, and winner determination is weighted set
packing on a universe of size $P$. The dynamic program
$DP[M]=\max\{DP[M\setminus T_b]+w_b:b\text{ uses }T_b\subseteq M\}$
over masks $M\subseteq[P]$ scans the bid list at each state, giving
$2^{O(P)}\cdot\mathrm{poly}(|\mathcal B|)$.

\section{Full Proofs for Section~\ref{sec:incentives}}

\begin{proof}[of Theorem~\ref{thm:incompat}]
Let there be one order. Solver $i$ has private deliverable value $v_i$
and may report any binding value $r_i\le v_i$; over-reporting is excluded
by settlement, while under-reporting is feasible. By non-wastefulness the
order is allocated whenever some report is positive, and by
value-respecting the winner has a maximal report. Feasibility implies the
winner can deliver at most its true value. Best-bid fairness requires the
delivery to be at least the endogenous benchmark $\max_i r_i$.

Let the highest and second-highest true values be
$v^{(1)}>v^{(2)}$. If the top solver reports truthfully, it sets the
benchmark to $v^{(1)}$, must deliver $v^{(1)}$, and earns zero surplus.
If instead it reports $r=v^{(2)}+\varepsilon$ with
$0<\varepsilon<v^{(1)}-v^{(2)}$, it still has the maximal report, the
benchmark falls to $v^{(2)}+\varepsilon$, and it can earn
$v^{(1)}-v^{(2)}-\varepsilon>0$. Truthful reporting is therefore not a
dominant strategy for any feasible non-wasteful value-respecting
mechanism enforcing the best-bid benchmark. The same argument shows that
any DSIC mechanism in this class can enforce at most the runner-up
benchmark.\hfill$\square$
\end{proof}

\begin{proof}[of Theorem~\ref{thm:posp}]
\emph{Endpoint.} If $g=1$, the floor and the batch value coincide, so
settling at the floor is optimal and strategyproof. Hence assume $g>1$
for the lower-bound expression.

\emph{Lower bound.} Restrict to the screening family $R_g$: a single
strategic solver holds one batch of private value $V\in[1,g]$ (floor mass
normalized to $1$, so $\FAIR(V)=V$), and the mechanism posts an outcome
as a function of the reported $V$. Write $p(V)$ for the probability the
batch is executed and $x(V)$ for the expected batch delivery (zero when
not executed). The fallback floor delivers welfare $1$, so
$W(V)=x(V)+1-p(V)$, and the solver's interim utility is
$u(V)=p(V)V-x(V)$. Because bids are binding, a type $V$ can only
\emph{under}report (claim $V'\le V$), so incentive compatibility implies
\[
u(V)\ge u(V')+(V-V')p(V') \qquad (V'\le V).
\]
Thus $u$ is nondecreasing and $u'(V)\ge p(V)$ a.e. Hence
$u(V)\ge\int_1^V p(t)\,dt$ and
\[
x(V)=p(V)V-u(V)\le
\psi(V):=p(V)V-\int_1^V p(t)\,dt .
\]
Integrating $\psi(V)/V^2$ over $[1,g]$ and applying Fubini gives
\[
\int_1^g\frac{\psi(V)}{V^2}\,dV
=\int_1^g\frac{p(V)}{V}\,dV
-\int_1^g p(t)\Big(\frac1t-\frac1g\Big)\,dt
=\frac1g\int_1^g p(t)\,dt
\le\frac{g-1}{g}.
\]
If the mechanism has competitive ratio $r$, i.e. $W(V)\ge V/r$
pointwise, then $x(V)\ge V/r-(1-p(V))\ge V/r-1$, and
\[
\int_1^g \frac{V/r-1}{V^2}\,dV
=\frac{\ln g}{r}-\Big(1-\frac1g\Big).
\]
Combining the two inequalities,
$(\ln g)/r\le2(1-1/g)$, so every randomized DSIC mechanism has
\[
r=\sup_{V\in[1,g]}\frac{V}{W(V)}
\ge\frac{g\ln g}{2(g-1)}
\ge\frac{\ln g}{2}.
\]

\emph{Posted-price upper bound.} Offer the batch a take-it-or-leave-it
delivered welfare on the geometric ladder $\{e^k:0\le k\le K\}$,
$K=\lceil\ln g\rceil$. Choose $k$ uniformly at random, execute the batch
if it accepts the posted threshold, and otherwise settle at the floors.
For any realized $V\in[e^k,e^{k+1})$, the draw $k$ (probability
$1/(K+1)$) executes the batch and delivers at least $e^k>V/e$, while all
draws deliver at least the floor mass. Hence
$\E[\text{welfare}]\ge\max\big(V/(e(K+1)),1\big)\ge\FAIR/(e(\lceil\ln
g\rceil+1))$.

\emph{Deterministic optimum.} Let the mechanism be deterministic. First,
its execution rule is a threshold: if $p(r)=1$ and $p(V)=0$ for some
$r<V$, type $V$ has truthful utility $0$, while under-reporting $r$ yields
utility $V-x(r)\ge V-r>0$ (feasibility at report $r$ gives $x(r)\le r$),
contradicting downward IC. Thus there is a threshold $\tau$ (up to the
boundary point) such that types below $\tau$ are not executed and types
above $\tau$ are. Second, the accepted delivery is bounded by $\tau$: for
$\tau\le r<V$, downward IC gives $V-x(V)\ge V-x(r)$, hence
$x(V)\le x(r)$; letting $r\downarrow\tau$ and using feasibility
$x(r)\le r$ gives $x(V)\le\tau$ for all executed types, in particular
$x(g)\le\tau$. Types just below $\tau$ have welfare $1$, while type $g$
has welfare at most $\tau$, so every deterministic DSIC rule has ratio at
least $\max\{\tau,g/\tau\}\ge\sqrt g$. The posted delivery $\sqrt g$
attains this bound. Since $R_g$ embeds as a degenerate single-minded
subfamily of the full model, all three bounds transfer.\hfill$\square$
\end{proof}

\begin{proof}[of Theorem~\ref{thm:lift}]
\emph{(i) Eligibility-posted packing.} In the public-floor regime, draw
$k\sim\mathrm{Unif}\{0,\dots,K\}$, $K=\lceil\ln g\rceil$, post the weight
$C_S e^{k}$ for each option $S$, and let a solver accept its
option iff its binding value clears this weight; among accepted options
take a maximum-weight disjoint packing using the posted weights. Accepted
options deliver their posted weights componentwise by the filter-passing
scaling condition from the mechanism model; all orders not covered by the
selected posted options settle at their public floors. The
posted weights and packing rule are report-independent, so each solver
faces a take-it-or-leave-it offer and the mechanism is universally
truthful. Let
$X$ be the total value of the batch options used by the offline fair
optimum. The geometric draw captures expected posted delivery at least
$X/(e(K+1))$. Independently, every realization settles all orders at least
at their public floors, so $W\ge C$. Since $\FAIR\le X+C$, the
max-inequality with $a=e(K+1)$ gives
$\E[W]\ge\FAIR/(e(\lceil\ln g\rceil+1)+1)$.

\emph{(ii) Coupled menu.} Draw $k$ uniformly from $\{0,\ldots,K\}$ and
set $t=e^k$. Post every option at price $C_S\,t$. Let the solver demand a
profit-maximizing disjoint collection $\mathcal C^\ast(t)$. This is
universally truthful because the posted menu is report-independent. It is
feasible because any demanded $S$ has $V_S\ge C_St\ge C_S$, and the
filter-passing scaling condition from the mechanism model implements total
delivery $C_St$ while clearing floors componentwise. Let
\[
p(t)=
\max_{\mathcal C\text{ disjoint}}
\sum_{S\in\mathcal C}(V_S-C_St).
\]
As a maximum of finitely many affine decreasing functions, $p$ is convex,
piecewise linear, nonincreasing, with $p(t)=0$ for $t\ge g$ (every
option has $V_S\le gC_S$). At any differentiability point
$-p'(t)=\sum_{S\in\mathcal C^\ast(t)}C_S$, so the delivered mass is
$D(t)=t\sum_{S\in\mathcal C^\ast(t)}C_S=t\,(-p'(t))$; at a kink the
right-derivative selects the continuing (smallest-floor) active
collection, and $D(t)\ge t\,(-p'_+(t))$. On $[t_k,t_{k+1}]=[e^k,e^{k+1}]$,
convexity gives $-p'$ nonincreasing, so
\[
\int_{t_k}^{t_{k+1}}(-p'(t))\,dt\le (-p'_+(t_k))(t_{k+1}-t_k)
=(-p'_+(t_k))(e-1)e^k\le (e-1)D(t_k).
\]
Summing over $k=0,\dots,K-1$ and using $p(e^K)=0$,
\[
p(1)=p(1)-p(e^K)=\int_1^{e^K}(-p'(t))\,dt\le(e-1)\sum_{k=0}^{K}D(e^k),
\]
so the uniform draw gives $\E[D]\ge p(1)/((e-1)(K+1))$. Since the fair
optimum's batches form a feasible disjoint collection at $t=1$,
$p(1)\ge\sum_{S\in B^\ast}(V_S-C_S)=\FAIR-C$, and combining with the
pointwise $W\ge C$ via the same $\max$-inequality yields
$\E[W]\ge\FAIR/((e-1)(\lceil\ln g\rceil+1)+1)$.

\emph{Independent geometric pricing fails.}
The separation uses the following construction. There are options
$S_t=\{0,t\}$ for $t=1,\dots,m$, all sharing the common order $0$, with
$c_0=c_t=\tfrac12$ and $C_{S_t}=1$. The solver has value $V_t=g$ for every
option. Let $P_t$ be the posted price of option $S_t$ and
$P_{\min}:=\min_t P_t$. Independent geometric prices make the solver
choose only the cheapest acceptable option: $D=P_{\min}$ if
$P_{\min}\le g$, and $D=0$
otherwise. Hence
\[
D\le1+g\,\mathbf 1\{P_{\min}\ge e\}.
\]
Moreover,
\[
\Pr[P_{\min}\ge e]=(K/(K+1))^m\le e^{-m/(K+1)} .
\]
Choosing $m=\lceil(K+1)\ln g\rceil$ gives $\E[D]\le2$. Thus
$\E[W]\le2+C=O(\log^2 g)$, while the fair optimum obtains
$\FAIR=g+(m-1)/2=\Omega(g)$ by taking one high-value option and settling
the remaining leaf orders at floors. Hence the independent geometric
analogue has ratio $\Omega(g/\log^2 g)$.\hfill$\square$
\end{proof}

\begin{proof}[of Theorem~\ref{thm:multi}]
\emph{Upper bound.} In random-priority coupled menus, condition on a
solver being served first (probability $1/N$): it faces the full coupled
menu on all orders, and the profit-curve telescope of
Theorem~\ref{thm:lift}(ii) applies verbatim, giving expected delivery at
least $p_i(1)/((e-1)(K+1))$ from that solver, where $p_i$ is computed on
the full order set. A solver served later faces a subset of orders and
contributes nonnegatively. Since each solver's optimal batches form a
disjoint collection on the full set, $\sum_i p_i(1)\ge\sum_{b\in
B^\ast}(V_b-C_b)=\FAIR-C$. Averaging over the uniform priority and
combining with $W\ge C$ gives
$\E[W]\ge\FAIR/(N(e-1)(\lceil\ln g\rceil+1)+1)$. Truthfulness holds
because a solver's menu depends only on which orders predecessors have
consumed, never on its own report.

\emph{Lower bound (symmetrized swarm).} For $g=1$ the lower bound is the
trivial constant bound, so assume $g>1$. Fix orders $o_1,\dots,o_N$ with
floors $\varepsilon=L/(N\rho)$ and auxiliary orders $u_1,\dots,u_L$ with
floor $1$; let $\mathcal U$ be the set of all orders, so
$C_{\mathcal U}\le2L$. Each identity
has the \emph{same} public option family
$\{\{o_1\},\ldots,\{o_N\},\mathcal U\}$; identities
differ only in private values, so any
report-independent service order is statistically independent of which
identity is the ``batcher.'' Choose $\rho\in(\max\{1,g/2\},g]$ avoiding
the (countably many) atoms of the $N$ multiplier distributions. The
batcher, placed by the adversary at an identity chosen after seeing the
mechanism, has $V(\mathcal U)=\rho\,C_{\mathcal U}$ and ratio-$1$
singletons. Every other identity (blocker $i$) has a ratio-$\rho$
singleton on $\{o_i\}$ and ratio-$1$ on all other options, including
$\mathcal U$. At multiplier $t$ a blocker is silent iff $t>\rho$;
write $q_i=\Pr[t_i>\rho]$. The batcher delivers
$\Theta(\rho C_{\mathcal U})$ only if
it is reached before any blocker has consumed an order of $\mathcal U$, i.e. only
if all predecessors are silent. Blockers together deliver at most
$N\varepsilon\rho=L$, and fallback floors contribute at most
$C_{\mathcal U}\le2L$,
so pointwise $W\le\mathrm{del}_{\mathrm{batcher}}+3L$. For any
realization of the order,
\[
\sum_{i}(1-q_i)\!\!\prod_{s\prec i}\!q_s \;=\; 1-\prod_s q_s \;\le\; 1,
\]
the $i$-th term being the probability that $i$ is the first non-silent
identity. Taking expectation over the (independent) order and the
multipliers,
\[
  \sum_i (1-q_i)\Pr[\text{predecessors of }i\text{ silent}] \le 1,
\]
so some identity $i^\ast$ has
$(1-q_{i^\ast})\Pr[\text{predecessors silent}]\le 1/N$. Placing the
batcher at $i^\ast$ caps its expected batch delivery at $2\rho L/N$,
while the fair optimum delivers at least $\rho L$ via the batch. Hence
\[
\frac{\FAIR}{\E[W]}\ge\frac{\rho}{3+2\rho/N}
=\Omega(\min(N,\rho))=\Omega(\min(N,g)),
\]
using $\rho>g/2$. The blockers' geometry is public and identical across
placements, so the bound holds even for geometry-aware orders.\hfill$\square$
\end{proof}

\section{Full Proofs for Section~\ref{sec:online}}

\begin{proof}[of Proposition~\ref{prop:online}]
Settle-at-deadline is fair and achieves $C$ on every instance, while
$\FAIR\le\OPT\le gC$, so the ratio is at most $g$; a staggered-deadline
instance forcing individual-only settlement shows it is at least $g$.\hfill$\square$
\end{proof}

\begin{proof}[of Theorem~\ref{thm:online}]
(i) If an almost-surely fair mechanism settles an exposed batch (a
covered order's deadline in the future) with positive probability on some
prefix, fix that prefix and append a higher standalone bid on a
later-deadline covered order: on the appended instance it has, with
positive probability, settled below a realized floor---not a.s.\ fair.
Hence exposed batches are never settled, and on the staggered hard family
the mechanism is individual-only a.s., with welfare $\le C$ against
$\FAIR=gC$.
(ii) Stopped welfare $q=V+C-C_S$ satisfies $q\in[C,gC]$ ($V\ge C_S$ and
$V\le gC_S$, $C_S\le C$). \emph{Deterministic} $\Theta(\sqrt g)$: accept
the first $q\ge\sqrt g\,C$; if $q^\ast\ge\sqrt g\,C$ the ratio is
$q^\ast/q\le\sqrt g$, else $W=C$ and ratio $<\sqrt g$; a two-step
adversary matches. \emph{Randomized} $O(1+\log g)$: with
$k\sim\mathrm{Unif}\{0,\dots,K\}$, accept the first $q\ge e^kC$. With
probability $1/(K+1)$, $e^kC\le q^\ast<e^{k+1}C$, and then the mechanism
accepts some offer with $q\ge e^kC>q^\ast/e$, so
$\E[W]\ge q^\ast/(e(K+1))$. The matching $\Omega(1+\log g)$ bound is
constant at $g=1$; for $g>1$ its logarithmic term is a Yao bound:
a persistent order with fleeting geometric batch offers $x_i=e^ic$,
$i\le\lfloor\ln g\rfloor$, shown for the first $M$ with
$\Pr[M\ge i]=e^{-i}$, caps every deterministic threshold at $2c$ while
$\E[\FAIR]=\Omega(c\log g)$. Greedy: a ratio-$g$ blocker arriving before a
ratio-$g$ batch is accepted (killing the batch) for any threshold
$T\le g$ and rejected with the batch for $T>g$, ratio $\Omega(g)$.\hfill$\square$
\end{proof}

\begin{proof}[of Theorem~\ref{thm:poa}]
In the per-order-decomposable individual game, the maximal standing bid
on $j$ both wins and delivers itself, so $c_j(b)=p_j(b)$. Let
$v_j^{\max}:=\max_i v_{i,j}$, choose
$i^\ast(j)\in\arg\max_i v_{i,j}$, and set
$\OPT_{\mathrm{ind}}:=\sum_j v_j^{\max}$. Let $\beta_j$ denote the
price/penalty term faced by the deviation on order $j$; below
$\beta_j\le c_j(b)=p_j(b)$. Let
$i^\ast(j)$ deviate to $W_j$ with density $1/(v^{\max}_j-w)$ on
$[0,(1-1/e)v^{\max}_j]$ (Syrgkanis--Tardos). Then the expected utility
from $j$ is $((1-\tfrac1e)v^{\max}_j-\beta_j)^+\ge(1-\tfrac1e)
v^{\max}_j-p_j(b)$ with $\beta_j\le c_j(b)=p_j(b)$, the penalty landing
under $b$. Summing and using
$\sum_iu_i(b)=\PW(b)-\sum_jp_j(b)$ in the CCE inequality cancels the
price terms, giving $\E[\PW]\ge(1-1/e)\OPT_{\mathrm{ind}}$; against the
unrestricted optimum, $\OPT\le g\,\OPT_{\mathrm{ind}}$, so the factor is
at most $(e/(e-1))g$.
The full-game obstruction is a latent batch present in $b_{-i}$
that executes under the deviation but not under $b$, breaking every
charging scheme.\hfill$\square$
\end{proof}

\section{Benchmark Manipulation}\label{app:manipulation}

On one directed pair, let an attacker of volume $v$ depress a quantile
benchmark by injecting mass $a$ at convex cost
$K(a)=\kappa_1a+\tfrac{\kappa_2}2a^2$, with local sensitivity $\beta$.
For a lower-tail quantile of $m$ honest quotes with density $f$ at
$b^\ast$, first-order sensitivity is $\beta=1/(mf(b^\ast))$.
The attacker chooses $a_t\in[0,\bar a]$ each round, with $\bar a>0$ and
$\kappa_2>0$. A public rate limit $r\in(0,1)$ bounds the downward bias
after a restart from $b_0=b^\ast$ by $b^\ast(1-(1-r)^t)$ in round $t$.
Steady-bias statements assume $\mu<1$.

\begin{theorem}[The contamination game]\label{thm:contam}
A marginal injection cost $\kappa_1\ge v\beta$ deters contamination.
Otherwise the optimal steady attack is
$a^\ast=\min\{\bar a,(v\beta-\kappa_1)/\kappa_2\}$, inducing relative
bias $\mu=\beta a^\ast/b^\ast$; bursts do not beat the constant attack,
and rate limits tax restarts through a finite ramp without changing the
steady state. Every public-floor guarantee is then read with
$G\mapsto G/(1-\mu)$.
\end{theorem}

\begin{proof}[of Theorem~\ref{thm:contam}]
The steady-state benchmark under a constant injection $a$ is
$b^\ast-\beta a$ (the rate limit does not bind once the bias is
constant), so the per-round profit is
$\pi(a)=v\beta a-\kappa_1 a-\tfrac{\kappa_2}{2}a^2$. With $\kappa_2>0$,
$\pi$ is strictly concave with $\pi'(0)=v\beta-\kappa_1$; hence if
$v\beta\le\kappa_1$ the unique maximizer is $a^\ast=0$ (no attack), and
otherwise $a^\ast=\min\{\bar a,(v\beta-\kappa_1)/\kappa_2\}$, with
relative bias $\mu=\beta a^\ast/b^\ast$.

\emph{Bursts do not help.} The published benchmark satisfies
$b_t\ge b^\ast-\beta a_t$ pointwise, so for every horizon $T$ and every
(possibly time-varying) policy $(a_t)$,
\[
\sum_{t=1}^T\big[v(b^\ast-b_t)-K(a_t)\big]
\le \sum_{t=1}^T\big[v\beta a_t-K(a_t)\big]
=\sum_{t=1}^T\pi(a_t)\le T\,\pi(a^\ast),
\]
and the constant policy $a^\ast$ attains $T\pi(a^\ast)-O(1)$ after the
ramp; hence no episodic schedule beats the constant attack.

\emph{Ramp tax (exact).} Starting from $b_0=b^\ast$, the rate limit caps
the attainable bias at round $t$ by $\Delta_t:=\min\{\Delta^\ast,\,
b^\ast(1-(1-r)^t)\}$ with $\Delta^\ast=\beta a^\ast$; reaching
$\Delta^\ast$ takes $T_\mu=\big\lceil\ln\tfrac{1}{1-\mu}\big/
\ln\tfrac{1}{1-r}\big\rceil$ rounds. An attacker who hugs the cap injects
$a_t=\Delta_t/\beta$, so the exact net rent forgone per restart is
\[
L_{\mathrm{net}}=\sum_{t=1}^{T_\mu}\Big[\pi(a^\ast)-\pi(\Delta_t/\beta)\Big]
\;\le\; \pi(a^\ast)\,T_\mu .
\]
When the cap is slack, so $a^\ast<\bar a$, and in the small-$(r,\mu)$
regime, $1-(1-r)^t\approx rt$ and the summand is
$\tfrac{\kappa_2}{2}(a^\ast)^2(1-rt/\mu)^2$, giving
$L_{\mathrm{net}}=\tfrac{\pi(a^\ast)\mu}{3r}\,(1+o(1))$; we state this only
as a first-order approximation, the exact figure being the finite sum
above.

\emph{Feedback.} A steady bias $\mu$ scales every floor on the pair by
$1-\mu$, so any guarantee whose benchmark input is the public-floor spread
$G=\max_{i,S:\,C_S>0}V_i(S)/C_S$ becomes $G/(1-\mu)$; the coverage law transforms
jointly,
$(G,\alpha)\mapsto(G/(1-\mu),\min\{1,\alpha/(1-\mu)\})$. Structural
production complementarity, directed-pair complexity, and produced-welfare
PoA are unaffected.\hfill$\square$
\end{proof}

\section{Deployment Diagnostics}\label{app:empirical}

\begin{center}
\footnotesize
\begin{minipage}{0.92\textwidth}
\centering
\textbf{Table 2.} Deployment diagnostics used only as motivation. The
visible spread is auction-level and observable; it is not the structural
complementarity factor $g$.
\medskip

\begin{tabular}{@{}ll@{}}
\toprule
Diagnostic & Value or interpretation\\
\midrule
Competitions in trace & 802,816 post-adoption solver competitions\\
Covered notional / gap & \$54.69B / \$5.13M (about 0.94 bps)\\
Visible spread tail & $\widehat G^{\mathrm{obs}}>2$: 25.3\%;
  $\widehat G^{\mathrm{obs}}>4$: 20.2\%\\
Small-notional tail & $\widehat G^{\mathrm{obs}}>2$: 63.1\%\\
Directed-pair support $P$ & directly observable from requested/paid pairs\\
Naive local reserve & 97.1\% vulnerable cell-windows in diagnostic\\
\bottomrule
\end{tabular}
\end{minipage}
\end{center}

\paragraph{Sample and filters.}
The trace contains 802,816 post-adoption solver competitions over an
eight-month window. The parser retains protocol-valid, non-filtered,
positive-score solver submissions. Full-sample quantities such as visible
spread and effective solver count use all parsed auctions. USD-denominated
quantities use smaller denominators because score conversion, notional
coverage, and sanity filters are available only for subsets of the trace.
These denominator changes are coverage restrictions, not inconsistencies.

\paragraph{Score variables.}
For each competition, $\mathrm{Score}_1$ is the winning solution score.
The reference score $\mathrm{Score}^{\mathrm{other}}_2$ is the best
positive score submitted by a solver different from the winner. The parser
groups by solver identity, removes all candidate solutions submitted by
the winning solver, and then takes the maximum remaining positive valid
score; if none exists, the reference score is zero. This best-non-winning
construction is solver-level: a winning solver's second candidate solution
is not a credible outside option.

\paragraph{Visible spread and accounting gap.}
When the denominator is positive,
$\widehat G^{\mathrm{obs}}=\mathrm{Score}_1/
\mathrm{Score}^{\mathrm{other}}_2$, with infinite spread when the
reference score is zero. The reported tail shares are computed on the full
802,816-auction sample. The dollar gap converts
$\mathrm{Score}_1-\mathrm{Score}^{\mathrm{other}}_2$ using ETH/USD prices,
winsorizes the dollar gap at the first and ninety-ninth percentiles, and
sums over the sanity-filtered denominator above. The \$54.69B covered
notional is the sell-token notional over the same denominator.

\paragraph{Reserve-vulnerability diagnostic.}
The benchmark diagnostic uses cell-windows from the reserve-construction
pipeline. With baseline quantile $\tau=0.25$, a cell-window is marked
vulnerable if the maximum eligible-sample share of any single solver is at
least $\tau$, meaning that a single solver could potentially influence the
local reserve quantile if its samples were contaminated. In the diagnostic,
97.1\% of auction-weighted cell-windows are vulnerable. The design
implication is to inherit from parent cells or use audit-only/soft reserves
when local support is sparse.


\end{document}
